\chapter*{Terminology}
\addcontentsline{toc}{chapter}{Terminology}

\begingroup
\small
\raggedright
\begin{description}[leftmargin=4.55cm,labelwidth=4.15cm,labelsep=0.4cm,style=sameline,itemsep=0.34em,topsep=0.35em,parsep=0pt]
  \item[Open Web Catcher (OWC)] The internal name of the complete evidence-collection platform built during the internship. It refers to the FastAPI backend, Next.js operator console, PostgreSQL persistence, MCP browser tooling, specialist-agent runtime, and review outputs used to catch open-web evidence from browser-visible streaming pages.
  \item[Agent] A model-driven runtime component that receives a bounded task, uses approved tools, and returns structured evidence or a failure reason.
  \item[Orchestrator] The control component that classifies the starting URL, prepares handoffs, schedules specialist agents, aggregates evidence, and decides the final run status.
  \item[Specialist agent] A browser-facing agent with a narrow responsibility, such as classification, landing-page discovery, hosting-page extraction, or embedded-player inspection.
  \item[Landing page] A page that lists events, channels, schedules, or match cards and is mainly useful for discovering downstream hosting candidates.
  \item[Hosting page] A page that usually owns play buttons, server/source controls, overlays, and redirects before a stream or embedded player becomes visible.
  \item[Embedded page] A direct player or iframe page used as a fallback extraction context when the hosting page exposes a player URL
  \item[Browser automation] Programmatic control of a real browser session to inspect rendered pages, click controls, manage frames, capture screenshots, and observe network activity.
  \item[Tool call] A browser, provider, memory, or runtime operation requested by an agent and recorded as part of the run trace.
  \item[Handoff] A short structured task brief prepared by the orchestrator for the next specialist agent, including target URL, evidence needed, route source, and navigation policy.
  \item[Trace] The ordered record of runtime events, model calls, tool calls, screenshots, status changes, costs, and final outputs for one investigation.
  \item[Observability] The ability to explain what happened during a run after it finishes, including which agent acted, which tools were used, and why a failure occurred.
  \item[Stream evidence] A URL, manifest, media request, screenshot, or player state that supports the claim that a target page exposed unauthorized streaming behavior.
  \item[Manifest] A playlist or media presentation file, such as HLS or MPEG-DASH, that points a player toward the media segments required for playback.
  \item[Iframe] An embedded browser frame that loads a player or third-party page inside the visible hosting page.
  \item[Provider enrichment] The step that resolves collected stream hosts or IP addresses into provider, country, network, or abuse-contact information.
  \item[Takedown email] A reviewable draft that groups stream evidence, provider details, screenshots, and rights-owner context for later human use.
  \item[Regression batch] A repeated set of website tests used to verify that a prompt, model, or tool change fixes new cases without breaking previously working ones.
  \item[Evidence artifact] Any persisted output used for review, including screenshots, streams, provider rows, emails, model usage, and tool results.
  \item[Runtime event] A timestamped event emitted during execution, such as agent start, tool success, screenshot capture, stream extraction, or failure.
  \item[Prompt] The role and task instruction compiled for an agent before it calls the language model.
  \item[Memory hint] A reusable site or layout note injected as soft guidance; it can help the agent but must be confirmed against live browser evidence.
  \item[Input tokens] Prompt, trace, tool-output, memory, and context text sent to the model during an LLM call.
  \item[Cached input tokens] Input tokens reused through the prompt or context cache instead of being billed and processed as fully new context.
  \item[Output tokens] Tokens generated by the model as an answer, tool decision, explanation, or structured result.
  \item[Total tokens] The aggregate token count used by the dashboard to summarize input and output usage for model calls.
  \item[Context window] The maximum amount of prompt, trace, tool output, and memory text that the model can consider in one call.
  \item[Prompt cache] A cache mechanism used to reuse repeated prompt and context content across similar model calls.
  \item[Cache hit rate] The share of model input served from cache within a specific dashboard or model-usage view. It must be read with its denominator.
  \item[LLM call] One recorded invocation of the language model, including the prompt, model name, token counters, and cost estimate.
  \item[Model spend] The recorded or estimated monetary cost of model calls for a run, batch, or dashboard snapshot.
  \item[Cost per run] Model spend normalized by the number of runs in the selected view; it is used to compare evaluation versions.
  \item[Tool success rate] Successful observed tool calls divided by observed tool calls in the same telemetry snapshot.
  \item[Stream yield] The share of runs that produced at least one stream, manifest, media, or embedded-player evidence candidate.
  \item[Strict stream success] A stricter success measure where a run counts only when it reaches reviewable stream evidence, not merely a completed status.
  \item[Pipeline run] One end-to-end investigation from a seed URL through classification, specialist routing, evidence collection, enrichment, and final status.
  \item[Active run] A pipeline run that is currently queued, running, or waiting on browser/model work.
  \item[Queued run] A run or task that has been accepted but has not yet started browser execution.
  \item[Run status] The final or current state of a pipeline run, such as success, partial, failed, cancelled, \texttt{no\_streams}, or \texttt{no\_hosting\_pages}.
  \item[\texttt{no\_streams}] A terminal status used when the browser path was explored but no reviewable media, manifest, stream, or embedded evidence was found.
  \item[\texttt{no\_hosting\_pages}] A terminal status used when a landing-page investigation did not discover credible downstream hosting candidates.
  \item[Evidence bundle] The grouped review output for one run, including streams, screenshots, provider rows, email drafts, status, costs, and trace links.
  \item[Provider coverage] The extent to which collected stream or host evidence could be enriched with provider, network, or registration context.
  \item[Dashboard snapshot] A live operator-console aggregate at a specific time. It should not be mixed blindly with controlled regression-batch metrics.
  \item[Persistence health] A dashboard view showing whether runs, agent runs, runtime events, model calls, tool calls, and outputs are being stored correctly.
  \item[Tokenized URL] A media or manifest URL containing temporary query parameters, signatures, or session-bound path values.
  \item[Queue message] A structured payload sent through a queue to start work or return results without blocking a synchronous web request.
  \item[\texttt{agent\_tasks}] The Azure Service Bus input queue used in the company deployment flow to receive target URLs and task metadata.
  \item[Results queue] The configured Azure Service Bus output queue used in the company deployment flow to return matched URLs, metadata, screenshots, rejected URLs, and failure information.
  \item[Browser profile] A browser context with its own state, cookies, cache, and execution settings for a run or agent.
  \item[Profile-scoped session] A browser session tied to one profile so page state can be preserved while avoiding uncontrolled sharing across unrelated runs.
  \item[Static collection] A monitoring step that fetches pages or links without completing the full rendered browser interaction path. In this project it remains useful for discovery but can miss evidence that appears only after scripts, clicks, frames, or playback.
  \item[Rendered browser state] The visible and runtime state produced after the browser executes JavaScript, handles frames, loads media, and reacts to user-like actions. OWC treats this state as an evidence surface.
  \item[Channel detection] The process of identifying a broadcaster or channel signal from landing hints, player labels, screenshots, OCR text, structured agent output, or stream metadata. Landing hints must be verified by hosting or embedded evidence when possible.
  \item[Fingerprint profile] A browser identity profile used to keep navigator, viewport, session, and related browser-surface behavior coherent during a run. It reduces obvious automation inconsistencies but does not guarantee bypass of anti-bot systems.
  \item[Proxy rotation] The runtime selection of a network proxy for browser execution. In OWC, proxy candidates are validated and kept sticky until failure to avoid breaking session-bound pages.
  \item[Sticky proxy] A proxy selected for a browser profile and reused while it remains healthy. This prevents identity changes in the middle of one investigation.
\end{description}
\endgroup
\clearpage
